\section{Run-edits}\label{sec:surgery}
\subsection{Insertion and deletion}
\begin{definition}\label{sync:surgery:2}\leavevmode
Let $()$ denote the empty tuple.
\begin{enumerate}[font=\itshape]
\item Let $U$ be a run. Let $i\in[0,|U|]_{\Z}$.
Let $v\in N(U)$. Let $j\ge1$ be an integer.
We refer to any one of the operations
\[
 U\mapsto\operatorname{Ins}(U,i,N),\qquad
 U\mapsto\operatorname{Del}(U,v),\qquad
 U\mapsto\operatorname{Ins}(U,i,R_j),
\]
or an exchange of two adjacent simple runs as a \textit{run-edit} of $U$.
\item We refer to a run-edit that inserts or deletes a vertical edge as a
\textit{vertical run-edit}.
\item Let $U_0$ be a run and $\ell\in\Z_{\ge0}$.
For each $j\in[\ell]_{\Z}$, let $\omega_j$ be a run-edit of $U_{j-1}$ and define
\[
 U_j=\omega_j(U_{j-1}).
\]
We call $\Omega=(\omega_j)_{j\in[\ell]_{\Z}}$ a \textit{run-edit list} from $U_0$ to $U_\ell$.
We define
\[
 \operatorname{RunEResult}(\Omega,U_0)=U_\ell,
 \qquad \operatorname{RunEResult}((),U_0)=U_0.
\]
\end{enumerate}
\end{definition}

\begin{lemma}[H-precursors after insertion]\label{lem:insert-precursor-chain}
Let $Y=Y_1\ccn\cdots\ccn Y_n$ be a run, let $g\in\operatorname{CtdG}(Y)$. Let $q\in[1,n-1]_{\Z}$. Put
\[
 h_0=Y_1,\qquad W=\operatorname{Ins}(Y,q,N),\qquad w_*=W_{q+1}.
\]
Define the map $\tau:\operatorname{NE}(Y)\to\operatorname{NE}(W)$ by
\[
 \tau(Y_k)=
 \begin{cases}
 W_k,&k\in[q]_{\Z},\\
 W_{k+1},&k\in[q+1,n]_{\Z}.
 \end{cases}
\]
Define $g'\in G(W)$ by
\begin{equation}\label{eq:insertion-grading}
 g'(w_*)=0,\qquad
 g'(\tau(h_0))=g(h_0)+1,\qquad
 g'(\tau(e))=g(e)\quad(e\in\operatorname{NE}(Y)\setminus\{h_0\}).
\end{equation}
Let
\[
 H_q=\{r\in[q]_{\Z}:Y_r\in E(Y),\
          \operatorname{Debt}_{Y,g}(Y_r,q)>0\}.
\]
The following statements hold.

\begin{enumerate}
\item\label{c:insert-active-chain} The set $H_q$ contains $1$. Its decreasing enumeration is uniquely determined:
\[
H_q=\{r_0,\ldots,r_d\},\qquad r_0>\cdots>r_d=1.
\]
For $j\in[0,d]_{\Z}$, put
\[
a_j=Y_{r_j},\qquad A_j=\operatorname{Node}_{Y,g}(a_j),\qquad
t_j=\operatorname{End}_{\mathrm{NE}}(A_j)=Y_{s_j}.
\]
Then, \begin{equation}\label{eq:insertion-node-chain}
A_0\subsetneq A_1\subsetneq\cdots\subsetneq A_d=Y,\qquad
q<s_0<\cdots<s_d=n,
\end{equation}
and $A_{j+1}$ is the parent of $A_j$ for $j\in[0,d-1]_{\Z}$.

\item\label{c:insert-chain-precursor} Use the indices from Claim~\ref{c:insert-active-chain}. The h-precursor $P_{W,g'}$ exists and satisfies
\begin{equation}\label{eq:insertion-chain-map}
\begin{aligned}
P_{W,g'}(w_*)&=\tau(a_0),\\
P_{W,g'}(\tau(t_j))&=\tau(a_{j+1})
   &&(j\in[0,d-1]_{\Z}),\\
P_{W,g'}(\tau(v))&=\tau(P_{Y,g}(v))
   &&\bigl(v\in N(Y)\setminus\{t_j:j\in[0,d-1]_{\Z}\}\bigr).
\end{aligned}
\end{equation}
For $d=0$, the middle line has no instances.

\item\label{c:insert-chain-fibers} For $e\in E(Y)$, put $F_e=P_{Y,g}^{-1}(\{e\})$. For a set of edges $F$, write $\tau(F)=\{\tau(v):v\in F\}$. Define
\[
b_0=w_*,\qquad b_j=\tau(t_{j-1})\quad(j\in[d]_{\Z}).
\]
Then, \begin{equation}\label{eq:insertion-chain-fibers}
\begin{aligned}
P_{W,g'}^{-1}(\{\tau(a_j)\})
   &=\tau(F_{a_j}\setminus\{t_j\})\sqcup\{b_j\}
     &&(j\in[0,d-1]_{\Z}),\\
P_{W,g'}^{-1}(\{\tau(a_d)\})
   &=\tau(F_{a_d})\sqcup\{b_d\},\\
P_{W,g'}^{-1}(\{\tau(e)\})
   &=\tau(F_e)
     &&\bigl(e\in E(Y)\setminus\{a_0,\ldots,a_d\}\bigr).
\end{aligned}
\end{equation}
Consequently, every $e\in E(Y)$ satisfies
\[
\#P_{W,g'}^{-1}(\{\tau(e)\})=g'(\tau(e)),
\]
and
\[
\operatorname{Node}_{W,g'}(\tau(h_0))=W.
\]
\end{enumerate}
\end{lemma}
\begin{proof}
\prf{Claim~\ref{c:insert-active-chain}} Let us enumerate $H_q$ and the nodes $\operatorname{Node}_{Y,g}(Y_r)$ for $r\in H_q$. Connectedness gives $Y=\operatorname{Node}_{Y,g}(h_0)$. Every positive-grade horizontal edge has a node by Claim~\ref{c:node-contained} of Lemma~\ref{lem:node-structure}. If $e=Y_r\in E(Y[0,q])$ and its node ends at $Y_s$, then
\begin{equation}\label{eq:insertion-active-end}
 \operatorname{Debt}_{Y,g}(e,q)>0\quad\Longleftrightarrow\quad q<s.
\end{equation}

Let us order these nodes by inclusion.
We have $P_{Y,g}(Y_n)=h_0$ and $q<n$.
Equation~\eqref{eq:insertion-active-end} gives $1\in H_q$.
Every node indexed by $H_q$ contains the boundary at $q$.
Claim~\ref{c:node-nesting} of Lemma~\ref{lem:node-structure} therefore orders them by inclusion.
The order $r_0>\cdots>r_d=1$ gives
\[
 A_0\subsetneq A_1\subsetneq\cdots\subsetneq A_d=Y.
\]
Claim~\ref{c:node-ends} of Lemma~\ref{lem:node-structure} gives
\[
 q<s_0<\cdots<s_d=n.
\]

Let us prove that $A_{j+1}$ is the parent of $A_j$ for $j\in[0,d-1]_{\Z}$.
If a node $\operatorname{Node}_{Y,g}(Y_r)$ satisfied
\[
 A_j\subsetneq\operatorname{Node}_{Y,g}(Y_r)\subsetneq A_{j+1},
\]
then it would contain the boundary at $q$. Hence,
\[
 r\in H_q,\qquad r_{j+1}<r<r_j.
\]
This contradicts the consecutive enumeration of $H_q$.
Definition~\ref{def:node-relations}(1) proves the parent assertion.

\prf{Claim~\ref{c:insert-chain-precursor}} Let us verify Definition~\ref{def:stack}(1)(a)--(c) for $Q$. Define $Q:N(W)\to E(W)$ by~\eqref{eq:insertion-chain-map}. Define its candidate debt, before knowing that $Q$ is an h-precursor, by
\begin{equation}\label{eq:insertion-candidate-debt}
 D_Q(x,\ell)=g'(x)-\#\{v\in N(W[0,\ell]):Q(v)=x\},
 \qquad x\in E(W[0,\ell]).
\end{equation}
Write $\mathbf1_{\mathcal P}=1$ when $\mathcal P$ holds and $0$ otherwise. For $k\in[0,q]_{\Z}$ and $e\in E(Y[0,k])$, equations~\eqref{eq:insertion-grading}, \eqref{eq:insertion-chain-map}, and~\eqref{eq:insertion-candidate-debt} give
\begin{equation}\label{eq:insertion-debt-before}
 D_Q(\tau(e),k)=\operatorname{Debt}_{Y,g}(e,k)+\mathbf1_{e=h_0}.
\end{equation}
For $k\in[q,n]_{\Z}$ and $e\in E(Y[0,k])$, equations~\eqref{eq:insertion-grading}, \eqref{eq:insertion-chain-map}, and~\eqref{eq:insertion-candidate-debt} give
\begin{equation}\label{eq:insertion-debt-after}
 \begin{aligned}
 D_Q(\tau(e),k+1)
 &=\operatorname{Debt}_{Y,g}(e,k)+\mathbf1_{e=h_0}-\mathbf1_{e=a_0}\\
 &\quad+\sum_{\substack{j\in[0,d-1]_{\Z}\\s_j\le k}}
       (\mathbf1_{e=a_j}-\mathbf1_{e=a_{j+1}}).
 \end{aligned}
\end{equation}
$w_*$ consumes one unit at $a_0$. Changing the owner of $t_j$ returns one unit at $a_j$ and consumes one unit at $a_{j+1}$.

Let us verify $Q(\tau(Y_k))=\tau(P_{Y,g}(Y_k))$ for $Y_k\in N(Y)$ with $k\le q$, and $Q(w_*)=\tau(a_0)$. The debt of $h_0$ is positive at every boundary before $n$. By~\eqref{eq:insertion-debt-before}, adding one unit to this first horizontal edge does not change the last horizontal edge of positive debt. Thus, $Q$ agrees with the h-precursor rule before $w_*$. At boundary $q$, the last positive-debt edge is $a_0$. Thus, $Q(w_*)=\tau(a_0)$.

Let us verify $Q(\tau(Y_k))$ for $Y_k\in N(Y)$ with $k>q$. Put $s_{-1}=q$. For $j\in[0,d]_{\Z}$, $s_{j-1}<k\le s_j$, and $e\in E(Y[0,k-1])$, the sum in~\eqref{eq:insertion-debt-after} telescopes to
\begin{equation}\label{eq:insertion-debt-between}
 D_Q(\tau(e),k)
 =\operatorname{Debt}_{Y,g}(e,k-1)+\mathbf1_{e=h_0}-\mathbf1_{e=a_j}.
\end{equation}
Here boundary $k$ of $W$ is immediately before $\tau(Y_k)=W_{k+1}$.

First, suppose $j<d$ and $Y_k$ is vertical with $k<s_j$. If $P_{Y,g}(Y_k)=a_j$, then both $Y_k$ and $t_j$ have owner $A_j$. Thus,
\begin{equation}\label{eq:insertion-two-units}
 \operatorname{Debt}_{Y,g}(a_j,k-1)\ge2,
 \qquad D_Q(\tau(a_j),k)\ge1.
\end{equation}
If $P_{Y,g}(Y_k)\ne a_j$, its h-precursor follows $a_j$, because $a_j$ has positive debt until $s_j$. Equation~\eqref{eq:insertion-debt-between} leaves its debt unchanged. By~\eqref{eq:insertion-debt-between} and~\eqref{eq:insertion-two-units}, in both cases $D_Q(\tau(P_{Y,g}(Y_k)),k)>0$ and $D_Q(\tau(Y_r),k)=0$ for every horizontal edge $Y_r$ strictly between $P_{Y,g}(Y_k)$ and $Y_k$. Hence, $Q(\tau(Y_k))=\tau(P_{Y,g}(Y_k))$.

Second, suppose $j<d$ and $k=s_j$. We have $t_j=Y_{s_j}$, $P_{Y,g}(t_j)=a_j$, and $s_j=\max\{k:Y_k\in N(Y),\ P_{Y,g}(Y_k)=a_j\}$. Its parent node ends at $s_{j+1}>s_j$ by~\eqref{eq:insertion-node-chain}. Hence, \eqref{eq:insertion-debt-between} gives
\begin{equation}\label{eq:insertion-ending-debts}
\begin{aligned}
 \operatorname{Debt}_{Y,g}(a_j,s_j-1)&=1,
 &D_Q(\tau(a_j),s_j)&=0,\\
 \operatorname{Debt}_{Y,g}(a_{j+1},s_j-1)&\ge1,
 &D_Q(\tau(a_{j+1}),s_j)&\ge1.
\end{aligned}
\end{equation}

Let us exclude a later horizontal edge of positive candidate debt when $j<d$ and $k=s_j$.
Let $r\in[r_{j+1}+1,s_j-1]_{\Z}\setminus\{r_j\}$ satisfy $Y_r\in E(Y)$.
If $g(Y_r)>0$, its node is contained in $A_j$, or is disjoint from $A_j$ and ends before $A_j$ starts.
Claim~\ref{c:insert-active-chain} excludes a node strictly between $A_j$ and $A_{j+1}$.
Claims~\ref{c:node-ends} and~\ref{c:node-nesting} of Lemma~\ref{lem:node-structure} give
\[
 \operatorname{Debt}_{Y,g}(Y_r,s_j-1)=0.
\]
If $g(Y_r)=0$, then $\operatorname{Debt}_{Y,g}(Y_r,s_j-1)=0$ by Claim~\ref{c:prefix-capacity} of Lemma~\ref{lem:precursor-prefix-counts} and Definition~\ref{def:stack}(3).
Equation~\eqref{eq:insertion-debt-between} therefore gives
\[
 D_Q(\tau(Y_r),s_j)=0.
\]
The equalities $D_Q(\tau(Y_r),s_j)=0$ for $r\in[r_{j+1}+1,s_j-1]_{\Z}\setminus\{r_j\}$ satisfying $Y_r\in E(Y)$, together with~\eqref{eq:insertion-ending-debts}, show that
$\tau(a_{j+1})$ is the last horizontal edge of positive candidate debt.
Thus, $Q(\tau(t_j))=\tau(a_{j+1})$.

Third, suppose $j=d$. The two indicator terms in~\eqref{eq:insertion-debt-between} cancel because $a_d=h_0$. Thus, $Q(\tau(Y_k))=\tau(P_{Y,g}(Y_k))$ for every $Y_k\in N(Y)$ with $s_{d-1}<k\le s_d$. Induction through the edges verifies all three conditions of Definition~\ref{def:stack}(1). Uniqueness in Claim~\ref{c:predecessor-uniqueness} of Lemma~\ref{lem:predecessor-uniqueness} gives $Q=P_{W,g'}$.

\prf{Claim~\ref{c:insert-chain-fibers}} Let us count the sets $P_{W,g'}^{-1}(\{\tau(e)\})$ for $e\in E(Y)$. For $j<d$, the fiber at $a_j$ loses $t_j$ and gains exactly $b_j$. The fiber at $a_d$ gains $b_d$ and loses no edge. Every other fiber is copied by $\tau$. These unions are disjoint: $w_*$ is new. For $j\in[d]_{\Z}$, we have $P_{Y,g}(t_{j-1})=a_{j-1}\ne a_j$. This proves~\eqref{eq:insertion-chain-fibers}. Connectedness of $g$ gives $\#F_e=g(e)$, including $F_e=\varnothing$ when $g(e)=0$. Therefore, $\#P_{W,g'}^{-1}(\{\tau(e)\})=g'(\tau(e))$ for every $e\in E(Y)$. The last edge of $P_{W,g'}^{-1}(\{\tau(h_0)\})$ is $\tau(Y_n)=W_{n+1}$. Thus, $\operatorname{Node}_{W,g'}(\tau(h_0))=W$. For $d=0$, this argument reads $P_{W,g'}^{-1}(\{\tau(h_0)\})=\tau(F_{h_0})\sqcup\{w_*\}$.
\end{proof}

\begin{lemma}[Balances after insertion]\label{lem:insert-balance-connected}
Use the hypotheses and construction of Lemma~\ref{lem:insert-precursor-chain}. Then, the following statements hold.

\begin{enumerate}
\item\label{c:insert-balances} For $\ell\in[0,n+1]_{\Z}$,
\[
\operatorname{bBal}_{W,g'}(\ell)=
\begin{cases}
\operatorname{bBal}_{Y,g}(\ell),&0\le\ell\le q,\\
\operatorname{bBal}_{Y,g}(\ell-1),&q+1\le\ell\le n+1.
\end{cases}
\]
Consequently,
\[
\operatorname{eBal}_{W,g'}(\tau(e))=\operatorname{eBal}_{Y,g}(e)
\quad(e\in\operatorname{NE}(Y)),\qquad
\operatorname{eBal}_{W,g'}(w_*)=\operatorname{bBal}_{Y,g}(q).
\]

\item\label{c:insert-connected-surplus} We have
\[
g'\in\operatorname{CtdG}(W),\qquad
\operatorname{Surp}(W,g')=\operatorname{Surp}(Y,g).
\]
If $m\in\Z_{\ge0}$, $g\in G_m(Y)$, and $g(h_0)+1\le m$, then $g'\in\operatorname{CtdG}_m(W)$.
\end{enumerate}
\end{lemma}
\begin{proof}
\prf{Claim~\ref{c:insert-balances}} Let us compute the two terms in the boundary balance.
Definition~\ref{def:colored-grading}(3)(a) gives
\[
 \operatorname{bBal}_{W,g'}(\ell)
 =\#E(W[0,\ell])-\sum_{v\in N(W[0,\ell])}g'(v)
 \qquad(\ell\in[0,n+1]_{\Z}).
\]

First, let us compute this expression for $\ell\in[0,q]_{\Z}$.
The insertion map $\tau$ gives
\[
 E(W[0,\ell])=\tau(E(Y[0,\ell])),\qquad
 N(W[0,\ell])=\tau(N(Y[0,\ell])).
\]
Equation~\eqref{eq:insertion-grading} gives $g'(\tau(v))=g(v)$ for every $v\in N(Y)$.
Hence,
\[
 \operatorname{bBal}_{W,g'}(\ell)
 =\#E(Y[0,\ell])-\sum_{v\in N(Y[0,\ell])}g(v)
 =\operatorname{bBal}_{Y,g}(\ell).
\]

Second, let us compute the expression for $\ell\in[q+1,n+1]_{\Z}$.
The prefix contains the new vertical edge $w_*$, and
\[
 \begin{aligned}
 E(W[0,\ell])&=\tau(E(Y[0,\ell-1])),\\
 N(W[0,\ell])&=\tau(N(Y[0,\ell-1]))\sqcup\{w_*\}.
 \end{aligned}
\]
Since $g'(w_*)=0$, we obtain
\[
 \operatorname{bBal}_{W,g'}(\ell)
 =\#E(Y[0,\ell-1])-\sum_{v\in N(Y[0,\ell-1])}g(v)-0
 =\operatorname{bBal}_{Y,g}(\ell-1).
\]

Let us evaluate the boundary balance after each edge.
For $k\in[q]_{\Z}$, we have $\tau(Y_k)=W_k$.
For $k\in[q+1,n]_{\Z}$, we have $\tau(Y_k)=W_{k+1}$.
Substituting $\operatorname{bBal}_{W,g'}(k)=\operatorname{bBal}_{Y,g}(k)$ for $k\le q$ and $\operatorname{bBal}_{W,g'}(k+1)=\operatorname{bBal}_{Y,g}(k)$ for $k>q$ gives
\[
 \operatorname{eBal}_{W,g'}(\tau(Y_k))
 =\operatorname{bBal}_{Y,g}(k)
 =\operatorname{eBal}_{Y,g}(Y_k)\qquad(k\in[n]_{\Z}).
\]
Since $w_*=W_{q+1}$, they also give
\[
 \operatorname{eBal}_{W,g'}(w_*)
 =\operatorname{bBal}_{W,g'}(q+1)=\operatorname{bBal}_{Y,g}(q).
\]

\prf{Claim~\ref{c:insert-connected-surplus}} Let us verify the node balance condition. Lemma~\ref{lem:insert-precursor-chain} gives the h-precursor, every positive-grade horizontal node, and the root node $W$. Let $v\in N(Y)$, write $a=P_{Y,g}(v)$. Let $b$ satisfy $P_{W,g'}(\tau(v))=\tau(b)$. Claim~\ref{c:insert-chain-precursor} of Lemma~\ref{lem:insert-precursor-chain} gives either $b=a$ or
\[
 \operatorname{Node}_{Y,g}(a)\subsetneq\operatorname{Node}_{Y,g}(b).
\]
The node balance condition for $g$, Claim~\ref{c:node-descendant-balances} of Lemma~\ref{lem:node-structure}, and Claim~\ref{c:insert-balances} give
\[
 \operatorname{eBal}_{W,g'}(\tau(v))
 =\operatorname{eBal}_{Y,g}(v)
 \ge\operatorname{eBal}_{Y,g}(a)
 \ge\operatorname{eBal}_{Y,g}(b)
 =\operatorname{eBal}_{W,g'}(\tau(b)).
\]
The inequality $\operatorname{eBal}_{Y,g}(a)\ge\operatorname{eBal}_{Y,g}(b)$ is equality when $a=b$.

Let us check the balance at the inserted edge.
For the inserted edge, Claim~\ref{c:insert-active-chain} of Lemma~\ref{lem:insert-precursor-chain} gives $r_0\le q<s_0$. Claim~\ref{c:node-boundary-balances} of Lemma~\ref{lem:node-structure} applied to $A_0=Y[r_0-1,s_0]$ yields
\[
 \operatorname{eBal}_{W,g'}(w_*)
 =\operatorname{bBal}_{Y,g}(q)
 \ge\operatorname{eBal}_{Y,g}(a_0)
 =\operatorname{eBal}_{W,g'}(P_{W,g'}(w_*)).
\]
Definition~\ref{def:connected-grading}(1) now gives connectedness.

Let us compute the surplus. We have
\[
 \#E(W)=\#E(Y),\qquad
 \sum_{v\in N(W)}g'(v)=\sum_{v\in N(Y)}g(v)+0.
\]
Substitution in Definition~\ref{def:connected-grading}(2) gives $\operatorname{Surp}(W,g')=\operatorname{Surp}(Y,g)$. If $g\in G_m(Y)$ and $g(h_0)+1\le m$, equation~\eqref{eq:insertion-grading} gives $g'(\tau(h_0))\le m$ and $g'(e)\le m$ for every $e\in\operatorname{NE}(W)\setminus\{\tau(h_0)\}$. Thus, $g'\in G_m(W)$.
\end{proof}

\begin{lemma}[Insertion]\label{lem:vertical-insert}
Let $M\in[m-2]_{\Z}$, $u\in[M]_{\Z}$, and $v\in[M]_{\Z}$.
Let $X$ be a run such that $\#N(A)\le M$ for every simple run $A\subseteq X$.
Let $c\in\Z_{\ge0}$.
Consider $(g,\gamma)\in\mathcal C_c(X)$.
Use the following construction.
\begin{hypotheses}
\item\label{hyp:insertion-shell} Let $Y=\operatorname{Shell}(X;u,v,0)$ with the shell grading $\widetilde g=\operatorname{GShell}_{g}(X;u,v,0)$.
\item\label{hyp:insertion-operation} Let $i\in[|X|]_{\Z}$.
Let
\[
 W=\operatorname{Ins}(Y,u+1+i,N)=R_u\ccn\operatorname{Ins}(X,i,N)\ccn R_v.
\]
Let $q=u+1+i$. Define $g'\in G(W)$ by
\[
 g'(W_k)=\begin{cases}
 u+2,&k=1,\\
 \widetilde g(Y_k),&k\in[2,q]_{\Z},\\
 0,&k=q+1,\\
 \widetilde g(Y_{k-1}),&k\in[q+2,|Y|+1]_{\Z}.
 \end{cases}
\]
\item\label{hyp:insertion-colors} Define
$\kappa\colon\operatorname{NE}(X)\hookrightarrow\operatorname{NE}(W)$ by
\begin{equation}\label{eq:insertion-edge-map}
 \kappa(X_k)=
 \begin{cases}
 W_{u+1+k},&k\in[i]_{\Z},\\
 W_{u+2+k},&k\in[i+1,|X|]_{\Z}.
 \end{cases}
\end{equation}
Let
\[
 \gamma'=\operatorname{Recol}(X,W,\kappa,g';g,\gamma).
\]
\end{hypotheses}
Then, we have the following.

\begin{enumerate}
\item\label{c:vertical-insert-connected} The grading $g'$ has the following properties.
\begin{enumerate}[label=(\alph*),ref=\theenumi(\alph*)]
\item\label{c:vertical-insert-connectivity} $g'\in\operatorname{CtdG}(W)$.
\item\label{c:vertical-insert-residual} We have
$\operatorname{Surp}(W,g')=\operatorname{Surp}(X,g)$.
\end{enumerate}
\item\label{c:vertical-insert-preserved} We have
\begin{equation}\label{eq:insertion-grade-color}
 g'(\kappa(X_k))=g(X_k),\qquad
 \gamma'(\kappa(X_k))=\gamma(X_k)
 \quad(k\in[|X|]_{\Z}).
\end{equation}
\item\label{c:vertical-insert-injective} Fix $X$, $i$, $u$, and $v$.
The construction~\ref{hyp:insertion-shell}--\ref{hyp:insertion-colors} gives an injective map
\[
 \Phi_c\colon\mathcal C_c(X)\longrightarrow\mathcal C_c(W)
\]
such that $\Phi_c(g,\gamma)=(g',\gamma')$ satisfies~\eqref{eq:insertion-grade-color}
for every $(g,\gamma)\in\mathcal C_c(X)$.
\end{enumerate}
\end{lemma}
\begin{proof}
\prf{Claim~\ref{c:vertical-insert-connectivity}} Let us apply the insertion lemmas to the shell grading.
Put $q=u+1+i$ and $h_0=\operatorname{Start}_{\mathrm{NE}}(Y)$.
The hypotheses give
\[
 1\le q<|Y|,\qquad \widetilde g(h_0)+1=u+2\le M+2\le m.
\]
Claim~\ref{c:shell-connected} of Lemma~\ref{lem:shell} gives
$\widetilde g\in\operatorname{CtdG}_m(Y)$.
Apply Lemma~\ref{lem:insert-precursor-chain} to $(Y,\widetilde g,q)$.
Claim~\ref{c:insert-connected-surplus} of Lemma~\ref{lem:insert-balance-connected}
gives $g'\in\operatorname{CtdG}_m(W)$.

\prf{Claim~\ref{c:vertical-insert-residual}}
Claim~\ref{c:insert-connected-surplus} of Lemma~\ref{lem:insert-balance-connected}
and Claim~\ref{c:shell-residual} of Lemma~\ref{lem:shell} give
\[
 \operatorname{Surp}(W,g')=\operatorname{Surp}(Y,\widetilde g)=\operatorname{Surp}(X,g).
\]

\prf{Claim~\ref{c:vertical-insert-preserved}} Let us verify~\eqref{eq:insertion-grade-color}.
By~\eqref{eq:insertion-edge-map}, the copy of $X_k$ in $W$ has index $u+1+k$ for $k\in[i]_{\Z}$ and $u+2+k$ for $k\in[i+1,|X|]_{\Z}$.
Assumptions~\ref{hyp:insertion-operation} and~\ref{hyp:insertion-colors} change no grade or color on these edges.
Thus, the map~\eqref{eq:insertion-edge-map} satisfies~\eqref{eq:insertion-grade-color}.

\prf{Claim~\ref{c:vertical-insert-injective}} Let us recover the input from the inserted grading.
Fix $X$, $i$, $u$, and $v$.
For $(g_1,\gamma_1)\in\mathcal C_c(X)$ and $(g_2,\gamma_2)\in\mathcal C_c(X)$, equation~\eqref{eq:insertion-grade-color} gives
\[
 \Phi_c(g_1,\gamma_1)=\Phi_c(g_2,\gamma_2)
 \ \Longrightarrow\ 
 (g_1(X_k),\gamma_1(X_k))=(g_2(X_k),\gamma_2(X_k))
 \quad(k\in[|X|]_{\Z}).
\]
Hence, $(g_1,\gamma_1)=(g_2,\gamma_2)$.
\end{proof}

\begin{lemma}[H-precursors after deletion]\label{lem:delete-precursors}
Let $Z$ be a tNE path. Let $g\in G(Z)$.
Suppose that the h-precursor $P_{Z,g}$ exists.
Let $v=Z_r\in N(Z)$ and put $x=P_{Z,g}(v)$.
Let $D=\operatorname{Del}(Z,v)$.
Define
\[
 \tau_v\colon\operatorname{NE}(Z)\setminus\{v\}\longrightarrow\operatorname{NE}(D),
 \qquad
 \tau_v(Z_k)=\begin{cases}D_k,&k<r,\\D_{k-1},&k>r.\end{cases}
\]
Define $g_v^-\in G(D)$ by
\[
 g_v^-(\tau_v(e))=
 \begin{cases}g(x)-1,&e=x,\\g(e),&e\ne x\end{cases}
 \qquad(e\in\operatorname{NE}(Z)\setminus\{v\}).
\]
Then, the h-precursor $P_{D,g_v^-}$ exists and satisfies
\begin{equation}\label{eq:delete-precursors}
 P_{D,g_v^-}(\tau_v(w))=\tau_v(P_{Z,g}(w))
 \qquad(w\in N(Z)\setminus\{v\}).
\end{equation}
\end{lemma}
\begin{proof}
Let us define $Q\colon N(D)\to E(D)$ by
\[
 Q(\tau_v(w))=\tau_v(P_{Z,g}(w))
 \qquad(w\in N(Z)\setminus\{v\}).
\]
The map $\tau_v$ preserves edge order and direction.
Condition (b) of Definition~\ref{def:stack}(1) gives $g(x)>0$.
Thus, $g_v^-$ has nonnegative grades.
For $Z_s\in N(Z)$ and $f\in E(Z[0,s-1])$, define
\[
 n_f(s)=\#\{w\in N(Z[0,s-1]):P_{Z,g}(w)=f\}.
\]
For $s\ne r$, define
\[
 n'_f(s)=\#\{w\in N(Z[0,s-1])\setminus\{v\}:Q(\tau_v(w))=\tau_v(f)\}.
\]
The definition of $Q$ gives
\[
 n'_f(s)=
 \begin{cases}
 n_f(s)-1,&s>r\text{ and }f=x,\\
 n_f(s),&s<r\text{ or }f\ne x.
 \end{cases}
\]

Let us verify that $Q$ satisfies the h-precursor conditions of Definition~\ref{def:stack}(1) at $\tau_v(Z_s)$ when $s<r$.
If $f=P_{Z,g}(Z_s)\ne x$, then $g_v^-(\tau_v(f))=g(f)$ and $n'_f(s)=n_f(s)$.
If $P_{Z,g}(Z_s)=x$, then
\[
 n_x(s)+1\le n_x(r)<g(x),\qquad n'_x(s)<g(x)-1.
\]
Thus, $Q(\tau_v(Z_s))$ satisfies conditions (a) and (b) of Definition~\ref{def:stack}(1).
For every $f\in E(Z[0,s-1])$, we have $n'_f(s)=n_f(s)$ and
$g_v^-(\tau_v(f))\le g(f)$.
Hence, every horizontal edge satisfying these two conditions after deletion
also satisfies them before deletion. Condition (c) follows.

Let us verify Definition~\ref{def:stack}(1)(a)--(c) for $Q$ at $\tau_v(Z_s)$ when $s>r$.
For $f\ne x$, we have $g_v^-(\tau_v(f))=g(f)$ and $n'_f(s)=n_f(s)$.
For $f=x$ and $s>r$, we have $v=Z_r\in N(Z[0,s-1])$ and $P_{Z,g}(v)=x$. Thus,
\[
 n_x(s)\ge1,\qquad
 n'_x(s)<g_v^-(\tau_v(x))
 \ \Longleftrightarrow\ n_x(s)-1<g(x)-1
 \ \Longleftrightarrow\ n_x(s)<g(x).
\]
Thus, the horizontal edges satisfying conditions (a) and (b) correspond under $\tau_v$.
The map $\tau_v$ preserves their order. Thus, condition (c) also holds.
Therefore, $Q$ is an h-precursor of $g_v^-$.
Claim~\ref{c:predecessor-uniqueness} of Lemma~\ref{lem:predecessor-uniqueness}
gives~\eqref{eq:delete-precursors}.
\end{proof}

\begin{lemma}[Deletion]\label{lem:vertical-delete}
Let $M\in[m-1]_{\Z}$.
Let $X$ be a run such that $\#N(A)\le M$ for every simple run $A\subseteq X$.
Let $v=X_r\in N(X)$. Let $c\in\Z_{\ge0}$.
Consider $(g,\gamma)\in\mathcal C_c(X)$.
Let $h=g(v)$. Let $u_0\in[M]_{\Z}$ and $v_0\in[M]_{\Z}$.
\begin{hypotheses}
\item\label{hyp:deletion-shell} Let $Y=\operatorname{Shell}(X;u_0,v_0,0)$ with the shell grading $\widetilde g=\operatorname{GShell}_{g}(X;u_0,v_0,0)$.
\item\label{hyp:deletion-edge} Let
\[
 W=\operatorname{Del}(Y,Y_{u_0+1+r})
 =R_{u_0}\ccn\operatorname{Del}(X,v)\ccn R_{v_0}.
\]
Let $q=u_0+1+r$. For each $x=X_j\in E(X[0,r-1])$ satisfying $g(x)>0$, define $g'_x\in G(W)$ by
\[
 g'_x(W_k)=\begin{cases}
 g(x)-1,&k=u_0+1+j,\\
 \widetilde g(Y_k),&k\in[1,q-1]_{\Z}\setminus\{u_0+1+j\},\\
 \widetilde g(Y_{k+1}),&k\in[q,|Y|-1]_{\Z}.
 \end{cases}
\]
\item\label{hyp:deletion-color} Define
$\kappa\colon\operatorname{NE}(X)\setminus\{v\}\hookrightarrow\operatorname{NE}(W)$ by
\[
 \kappa(X_k)=
 \begin{cases}
 W_{u_0+1+k},&k\in[r-1]_{\Z},\\
 W_{u_0+k},&k\in[r+1,|X|]_{\Z}.
 \end{cases}
\]
For each $x\in E(X[0,r-1])$ satisfying $g(x)>0$, define
\[
 \gamma'_x=\operatorname{Recol}(X,W,\kappa,g'_x;g,\gamma).
\]
\end{hypotheses}
Then, we have the following.

\begin{enumerate}
\item\label{c:vertical-delete-connected}\label{c:vertical-delete-connectivity}
There exists $y\in E(X[0,r-1])$ with $g(y)>0$ such that
$g'_y\in\operatorname{CtdG}(W)$.
\item\label{c:vertical-delete-residual} For every $x\in E(X[0,r-1])$ with $g(x)>0$, we have
\[
 \#E(W)-1-\sum_{w\in N(W)}g'_x(w)=\operatorname{Surp}(X,g)+h.
\]
\item\label{c:vertical-delete-preserved} For every $x\in E(X[0,r-1])$ with $g(x)>0$, we have
\begin{equation}\label{eq:deletion-grade-color}
\begin{aligned}
 g'_x(\kappa(x))&=g(x)-1,\\
 g'_x(\kappa(e))&=g(e),\qquad \gamma'_x(\kappa(e))=\gamma(e)
 \quad(e\in\operatorname{NE}(X)\setminus\{x,v\}).
\end{aligned}
\end{equation}
\item\label{c:vertical-delete-map} Fix $X$, $v$, $u_0$, and $v_0$.
Also, for each $(g,\gamma)\in\mathcal C_c(X)$, fix an edge $y(g,\gamma)$ satisfying Claim~\ref{c:vertical-delete-connected}.
The construction~\ref{hyp:deletion-shell}--\ref{hyp:deletion-color} gives a map
\[
 \Phi_c\colon\mathcal C_c(X)\longrightarrow
 \bigcup_{s\in[c,c+m]_{\Z}}\mathcal C_s(W)
\]
such that $\Phi_c(g,\gamma)=(g'_{y(g,\gamma)},\gamma'_{y(g,\gamma)})$.
The identities~\eqref{eq:deletion-grade-color} hold at $x=y(g,\gamma)$.
Also, we have
\[
 \operatorname{Surp}(W,g'_{y(g,\gamma)})=c+g(v).
\]
\end{enumerate}
\end{lemma}
\begin{proof}
\prf{Claim~\ref{c:vertical-delete-connectivity}} Let us choose $y\in E(X[0,r-1])$.
Let $y=P_{X,g}(v)$.
Definition~\ref{def:stack}(1) gives $y\in E(X[0,r-1])$ and $g(y)>0$.
Let $q=u_0+1+r$ and let $\widetilde y$ be the copy of $y$ in $Y$.
Claim~\ref{c:shell-connected} of Lemma~\ref{lem:shell} gives
$\widetilde g\in\operatorname{CtdG}_m(Y)$.
Claim~\ref{c:shell-copy-precursor} of Lemma~\ref{lem:shell} gives
$P_{Y,\widetilde g}(Y_q)=\widetilde y$.

Let us apply the deletion formula to the copied edge $Y_q$.
Let $\tau\colon\operatorname{NE}(Y)\setminus\{Y_q\}\to\operatorname{NE}(W)$
be the map $\tau_{Y_q}$ of Lemma~\ref{lem:delete-precursors}.
That lemma applies to $Y$, $\widetilde g$, and $Y_q$ and gives
\begin{equation}\label{eq:vertical-delete-precursor}
 P_{W,g'_y}(\tau(w))=\tau(P_{Y,\widetilde g}(w))
 \qquad(w\in N(Y)\setminus\{Y_q\}).
\end{equation}

Let us verify the node balance condition.
For $k\in[|Y|]_{\Z}\setminus\{q\}$, deletion of the vertical grade $h$ gives
\begin{equation}\label{eq:vertical-delete-balance}
 \operatorname{eBal}_{W,g'_y}(\tau(Y_k))
 =\operatorname{eBal}_{Y,\widetilde g}(Y_k)
 +\begin{cases}0,&k<q,\\h,&k>q.\end{cases}
\end{equation}
Let $w=Y_k\in N(Y)\setminus\{Y_q\}$ and let
$P_{Y,\widetilde g}(w)=Y_j$.
We have $j<k$. Thus,~\eqref{eq:vertical-delete-balance} gives
\[
\begin{aligned}
 &\operatorname{eBal}_{W,g'_y}(\tau(w))
 -\operatorname{eBal}_{W,g'_y}(\tau(Y_j))\\
 &\quad\ge\operatorname{eBal}_{Y,\widetilde g}(w)
 -\operatorname{eBal}_{Y,\widetilde g}(Y_j)\ge0.
\end{aligned}
\]
The inequality $\operatorname{eBal}_{Y,\widetilde g}(w)-\operatorname{eBal}_{Y,\widetilde g}(Y_j)\ge0$ follows from connectedness of $\widetilde g$ and
Claim~\ref{c:node-contained} of Lemma~\ref{lem:node-structure}.
Together with~\eqref{eq:vertical-delete-precursor}, this proves the node balance condition for $g'_y$.

Let us verify $\operatorname{Node}_{W,g'_y}(W_1)=W$.
The edges $Y_1$ and $Y_{|Y|}$ belong to the copies of $R_{u_0}$ and $R_{v_0}$, respectively.
Deleting $Y_q$ keeps $Y_1$ and $Y_{|Y|}$ and gives $g'_y(W_1)=\widetilde g(Y_1)$.
We have $P_{Y,\widetilde g}(Y_q)=\widetilde y\ne Y_1$.
Equation~\eqref{eq:vertical-delete-precursor} therefore gives $P_{W,g'_y}(\tau(w))=W_1$ for every $w\in P_{Y,\widetilde g}^{-1}(\{Y_1\})$.
Since $\operatorname{Node}_{Y,\widetilde g}(Y_1)=Y$, we obtain
\[
 \operatorname{Node}_{W,g'_y}(W_1)=W.
\]
Definition~\ref{def:connected-grading}(1) gives $g'_y\in\operatorname{CtdG}(W)$.
The grading $g'_y$ decreases only the positive grade $g(y)$.
Since $\widetilde g\in G_m(Y)$, we also have $g'_y\in G_m(W)$.

\prf{Claim~\ref{c:vertical-delete-residual}} Let us sum the vertical grades.
Let $x\in E(X[0,r-1])$ satisfy $g(x)>0$.
Deletion preserves the number of horizontal edges and removes vertical grade $h$.
By~\eqref{eq:residual} and Claim~\ref{c:shell-residual} of Lemma~\ref{lem:shell},
\[
 \#E(W)-1-\sum_{w\in N(W)}g'_x(w)
 =\#E(Y)-1-\left(\sum_{w\in N(Y)}\widetilde g(w)-h\right)
 =\operatorname{Surp}(X,g)+h.
\]

\prf{Claim~\ref{c:vertical-delete-preserved}} Let us identify the unchanged grades and colors.
Let $x\in E(X[0,r-1])$ with $g(x)>0$.
Assumption~\ref{hyp:deletion-edge} gives the grade identities
in~\eqref{eq:deletion-grade-color}.
For $e\in\operatorname{NE}(X)\setminus\{x,v\}$, we have
$g'_x(\kappa(e))=g(e)$.
Definition~\ref{def:recolor} and assumption~\ref{hyp:deletion-color}
therefore give $\gamma'_x(\kappa(e))=\gamma(e)$.

\prf{Claim~\ref{c:vertical-delete-map}} Let us prove
$\Phi_c(g,\gamma)\in\bigcup_{s\in[c,c+m]_{\Z}}\mathcal C_s(W)$.
Let $y=y(g,\gamma)$.
For the chosen $y$, Claim~\ref{c:vertical-delete-connectivity} gives
$g'_y\in\operatorname{CtdG}(W)$.
Definition~\ref{def:recolor} gives $\gamma'_y\in\operatorname{Col}(W,g'_y,m)$.
Definition~\ref{def:connected-grading}(2), Claim~\ref{c:vertical-delete-residual}, and $g(v)\in[0,m]_{\Z}$ give
\[
 c\le\operatorname{Surp}(W,g'_y)=c+g(v)\le c+m.
\]
Claim~\ref{c:vertical-delete-preserved} gives~\eqref{eq:deletion-grade-color} at $x=y$.
\end{proof}

\begin{lemma}[Insertion of a simple run]\label{lem:simple-run-insertion}
Let $Q$ be a tNE path, $c\in\Z_{\ge0}$, $q\in[1,|Q|-1]_{\Z}$, and $j\in[m]_{\Z}$.
Let $Q'=\operatorname{Ins}(Q,q,R_j)$.
Define $\tau\colon\operatorname{NE}(Q)\hookrightarrow\operatorname{NE}(Q')$ by
\begin{equation}\label{eq:run-insertion-map}
 \tau(Q_k)=
 \begin{cases}
 Q'_k,&k\in[q]_{\Z},\\
 Q'_{k+j+1},&k\in[q+1,|Q|]_{\Z}.
 \end{cases}
\end{equation}
Let $a_*=Q'_{q+1}$ and $v_t=Q'_{q+1+t}$ for $t\in[j]_{\Z}$.
For $(h,\eta)\in\mathcal C_c(Q)$, define $h'\in G_m(Q')$ by
\begin{equation}\label{eq:run-insertion-grades}
 h'=h|_{Q[0,q]}\ccn g_j^{\mathrm{run}}\ccn h|_{Q[q,|Q|]}.
\end{equation}
Let
\begin{equation}\label{eq:run-insertion-recolor}
 \eta'=\operatorname{Recol}(Q,Q',\tau,h';h,\eta).
\end{equation}
Then, the map
\[
 \Phi_{q,j,c}\colon\mathcal C_c(Q)\longrightarrow\mathcal C_{c+1}(Q'),
 \qquad \Phi_{q,j,c}(h,\eta)=(h',\eta')
\]
is well-defined and injective.
\end{lemma}
\begin{proof}
Let us determine the h-precursor after insertion of the simple run.
The new horizontal edge $a_*$ is the h-precursor of the $j$ new vertical edges.
The debt at $a_*$ decreases from $j$ to zero before the next original edge.
Definition~\ref{def:stack}(1) gives
\begin{equation}\label{eq:run-insertion-precursors}
\begin{aligned}
 P_{Q',h'}(v_t)&=a_*&&(t\in[j]_{\Z}),\\
 P_{Q',h'}(\tau(w))&=\tau(P_{Q,h}(w))&&(w\in N(Q)),\\
 \operatorname{Node}_{Q',h'}(a_*)&=Q'[q,q+j+1].
\end{aligned}
\end{equation}
For each $e\in E(Q)$, equation~\eqref{eq:run-insertion-precursors} gives
\[
 \#P_{Q',h'}^{-1}(\{\tau(e)\})=\#P_{Q,h}^{-1}(\{e\})=h(e)=h'(\tau(e)).
\]
For $a_0=\operatorname{Start}_{\mathrm{NE}}(Q)$, we therefore have
\[
 \operatorname{Node}_{Q',h'}(\tau(a_0))=Q'.
\]

Let us verify the node balance condition.
For $w=Q_k$, let $\varepsilon(w)=0$ when $k\le q$ and $\varepsilon(w)=1$ when $k>q$.
Equation~\eqref{eq:run-insertion-grades} gives
\begin{equation}\label{eq:run-insertion-balances}
\begin{aligned}
 \operatorname{eBal}_{Q',h'}(\tau(w))&=\operatorname{eBal}_{Q,h}(w)+\varepsilon(w),\\
 \operatorname{eBal}_{Q',h'}(a_*)&=\operatorname{eBal}_{Q',h'}(v_t)
 =\operatorname{bBal}_{Q,h}(q)+1\quad(t\in[j]_{\Z}).
\end{aligned}
\end{equation}
For $w\in N(Q)$ and $a=P_{Q,h}(w)$, the edge $a$ precedes $w$.
Hence, $\varepsilon(a)\le\varepsilon(w)$ and
\[
\begin{aligned}
 &\operatorname{eBal}_{Q',h'}(\tau(w))-\operatorname{eBal}_{Q',h'}(\tau(a))\\
 &\quad=\operatorname{eBal}_{Q,h}(w)-\operatorname{eBal}_{Q,h}(a)
       +\varepsilon(w)-\varepsilon(a)\ge0.
\end{aligned}
\]
Together with~\eqref{eq:run-insertion-precursors} and~\eqref{eq:run-insertion-balances}, this proves that $h'$ is connected.
Finally,
\begin{equation}\label{eq:run-insertion-surplus}
 |Q'|_E=|Q|_E+1,\qquad
 \operatorname{TNG}_{Q',h'}=\operatorname{TNG}_{Q,h},\qquad
 \operatorname{Surp}(Q',h')=\operatorname{Surp}(Q,h)+1.
\end{equation}

Let us recover $(h,\eta)$ from $\Phi_{q,j,c}(h,\eta)$.
Definition~\ref{def:recolor} and~\eqref{eq:run-insertion-grades} give
\[
 (h'(\tau(e)),\eta'(\tau(e)))=(h(e),\eta(e))
 \qquad(e\in\operatorname{NE}(Q)).
\]
Restriction to the image of $\tau$ therefore recovers $(h,\eta)$.
Hence, $\Phi_{q,j,c}$ is injective.
\end{proof}

\subsection{Composition of run-edits}
\begin{lemma}[Surplus reduction]\label{lem:surplus-reduction}
Let $m\in\Z_{\ge3}$ and $n\in\Z_{\ge0}$.
Let $Q$ be a tNE path.
Fix $\alpha=(a_j)_{j\in[n]_{\Z}}$ and
$\beta=(b_j)_{j\in[n]_{\Z}}$ with $a_j\in[m-1]_{\Z}$ and $b_j\in[m-1]_{\Z}$.
Define
\[
 Q_0=Q,\qquad Q_j=R_{a_j}\ccn Q_{j-1}\ccn R_{b_j}\quad(j\in[n]_{\Z}),
 \qquad \rho=\sum_{j=1}^n(a_j+1).
\]
Then, there is an injective map
\begin{equation}\label{eq:surplus-reduction-map}
 \Delta_{Q,\alpha,\beta}\colon
 \bigcup_{c\in[0,n]_{\Z}}\mathcal C_c(Q)\longrightarrow\mathcal C_0(Q_n)
\end{equation}
such that for each $(g,\gamma)\in\bigcup_{c\in[0,n]_{\Z}}\mathcal C_c(Q)$,
$(g',\gamma')=\Delta_{Q,\alpha,\beta}(g,\gamma)$ satisfies
\begin{equation}\label{eq:surplus-reduction-restriction}
 (g'((Q_n)_{\rho+k}),\gamma'((Q_n)_{\rho+k}))
 =(g(Q_k),\gamma(Q_k))\qquad(k\in[|Q|]_{\Z}).
\end{equation}
\end{lemma}
\begin{proof}
Let us construct the shell gradings and colorings.
Let $(g,\gamma)\in\bigcup_{c\in[0,n]_{\Z}}\mathcal C_c(Q)$.
Let $c=\operatorname{Surp}(Q,g)$ and $(g_0,\gamma_0)=(g,\gamma)$.
For $j\in[n]_{\Z}$, define
\[
 \epsilon_j(c)=\begin{cases}1,&j\le c,\\0,&j>c,\end{cases}
 \qquad
 g_j=\operatorname{GShell}_{g_{j-1}}(Q_{j-1};a_j,b_j,\epsilon_j(c)).
\]
Define $\iota_j\colon\operatorname{NE}(Q_{j-1})\hookrightarrow\operatorname{NE}(Q_j)$ by
\[
 \iota_j((Q_{j-1})_k)=(Q_j)_{a_j+1+k}
 \qquad(k\in[|Q_{j-1}|]_{\Z}).
\]
Let
\[
 \gamma_j=\operatorname{Recol}
 (Q_{j-1},Q_j,\iota_j,g_j;g_{j-1},\gamma_{j-1}).
\]

Let us prove that the final surplus is zero.
Claims~\ref{c:shell-connected}--\ref{c:shell-recolor} of Lemma~\ref{lem:shell} give, by induction,
\[
 (g_j,\gamma_j)\in\mathcal C_{\max(c-j,0)}(Q_j)
 \qquad(j\in[0,n]_{\Z}).
\]
For $j\le c$, we have $\epsilon_j(c)=1$ and
$\operatorname{Surp}(Q_{j-1},g_{j-1})=c-j+1\ge1$.
Thus, condition~\ref{hyp:shell-draining} of Lemma~\ref{lem:shell} applies.
For $j>c$, we have $\epsilon_j(c)=0$ and use condition~\ref{hyp:shell-preserving} of Lemma~\ref{lem:shell}.
Claim~\ref{c:shell-residual} of Lemma~\ref{lem:shell} gives
$\operatorname{Surp}(Q_j,g_j)=\operatorname{Surp}(Q_{j-1},g_{j-1})-\epsilon_j(c)$.
Since $c\le n$, define
\[
 \Delta_{Q,\alpha,\beta}(g,\gamma)=(g_n,\gamma_n)
 \in\mathcal C_0(Q_n).
\]

Let us recover $(g,\gamma)$ from $(g_n,\gamma_n)$.
Claim~\ref{c:shell-recolor} of Lemma~\ref{lem:shell} gives
\begin{equation}\label{eq:surplus-reduction-step-restriction}
 (g_j(\iota_j(e)),\gamma_j(\iota_j(e)))
 =(g_{j-1}(e),\gamma_{j-1}(e))
 \qquad(e\in\operatorname{NE}(Q_{j-1})).
\end{equation}
Composing~\eqref{eq:surplus-reduction-step-restriction} for $j\in[n]_{\Z}$ gives~\eqref{eq:surplus-reduction-restriction}.
Restriction to the copy of $Q$ therefore recovers $(g,\gamma)$. Thus,
$\Delta_{Q,\alpha,\beta}$ is injective on its whole domain.
\end{proof}

\begin{definition}\label{def:fiber-bound}
For $\nu\in\Z_{\ge0}$ and $D\in\Z_{\ge0}$, define
\[
 \operatorname{Fib}_m(\nu,D)=\binom{\nu+D}{D}2^{2mD}.
\]
\end{definition}

\begin{proposition}\label{prop:composition}
Let $M\in[m-1]_{\Z}$.
Let $U$ be a run with $\nu$ simple runs.
Let $K\in\Z_{\ge0}$, $A\in\Z_{\ge0}$, $D\in\Z_{\ge0}$, and $\ell=K+A+D$.
Assume the following.
\begin{hypotheses}
\item\label{hyp:comparison-edit-list} Let $\Omega=(\omega_s)_{s\in[\ell]_{\Z}}$ be a run-edit list starting at $U$.
Let
\[
 U_s=\operatorname{RunEResult}((\omega_1,\ldots,\omega_s),U)
 \quad(s\in[0,\ell]_{\Z}),\qquad V=\operatorname{RunEResult}(\Omega,U).
\]
Let $I_R$, $I_+$, and $I_-$ satisfy
\[
 [\ell]_{\Z}=I_R\sqcup I_+\sqcup I_-,\qquad \#I_R=K,\qquad \#I_+=A,\qquad \#I_-=D.
\]
For $s\in I_R\cup I_+$, fix $i_s\in[0,|U_{s-1}|]_{\Z}$.
For $s\in I_R$, fix $j_s\in\Z_{>0}$.
For $s\in I_-$, fix $v_s\in N(U_{s-1})$.
The run-edits satisfy
\[
 U_s=
 \begin{cases}
 \operatorname{Ins}(U_{s-1},i_s,R_{j_s}),&s\in I_R,\\
 \operatorname{Ins}(U_{s-1},i_s,N),&s\in I_+,\\
 \operatorname{Del}(U_{s-1},v_s),&s\in I_-.
 \end{cases}
\]
\item\label{hyp:comparison-run-lengths} For $s\in[0,\ell]_{\Z}$, the path $U_s$ is a run
and every simple run $B\subseteq U_s$ satisfies $\#N(B)\in[M]_{\Z}$.
\item\label{hyp:comparison-insertion-range} If $K+A>0$, then
$m\ge4$ and $M\le m-2$.
\item\label{hyp:comparison-shell-budget} Let $u\in[M]_{\Z}$ and $v\in[M]_{\Z}$.
Let $P$ be an integer satisfying
\begin{equation}\label{eq:layer-budget}
 P\ge K+A+(m+1)D.
\end{equation}
\item\label{hyp:comparison-target} Fix the tNE paths $W_{\mathrm L}$ and $W_{\mathrm R}$. Each path consists of $P$ simple runs.
Every simple run of these paths has vertical length in $[M]_{\Z}$.
Let
\begin{equation}\label{eq:target-shape}
 V_*=W_{\mathrm L}\ccn R_u\ccn V\ccn R_v\ccn W_{\mathrm R}.
\end{equation}
\end{hypotheses}
Then, we have the following.

\begin{enumerate}
\item\label{c:composition-target} There is a map
\[
 \Phi\colon\mathcal C_0(U)\longrightarrow\mathcal C_0(V_*)
\]
such that
\[
 \#\Phi^{-1}(\{\xi\})\le \operatorname{Fib}_m(\nu,D)
 \qquad(\xi\in\mathcal C_0(V_*)).
\]
\item\label{c:composition-bound} We have
\begin{equation}\label{eq:composition}
 C_0(V_*)\ge\frac{C_0(U)}{\operatorname{Fib}_m(\nu,D)}.
\end{equation}
\end{enumerate}
\end{proposition}
\begin{proof}
\prf{Claim~\ref{c:composition-target}} Let us construct $\Phi$.
Let $(g,\gamma)\in\mathcal C_0(U)$.
For $u$ and $v$ in assumption~\ref{hyp:comparison-shell-budget}, let
\[
 Y=\operatorname{Shell}(U;u,v,0)
 =R_u\ccn U\ccn R_v.
\]
Let $h=\operatorname{GShell}_{g}(U;u,v,0)$.
Define $\iota\colon\operatorname{NE}(U)\hookrightarrow\operatorname{NE}(Y)$ by
$\iota(U_k)=Y_{u+1+k}$ for $k\in[|U|]_{\Z}$ and let
\[
 \eta=\operatorname{Recol}(U,Y,\iota,h;g,\gamma).
\]
Claims~\ref{c:shell-connected} and~\ref{c:shell-residual} of Lemma~\ref{lem:shell} give $(h,\eta)\in\mathcal C_0(Y)$.
Suppose $P=0$.
Inequality~\eqref{eq:layer-budget} gives $K=A=D=\ell=0$.
Thus, $V=U$ and $W_{\mathrm L}=W_{\mathrm R}=\emptyset$.
The paths $Y$ and $V_*$ are nonempty translated copies.
Let
\[
 \mathbf t=\operatorname{Start}_{\mathrm{pt}}(V_*)-
             \operatorname{Start}_{\mathrm{pt}}(Y),\qquad
 \Phi(g,\gamma)=(h+\mathbf t,\eta+\mathbf t).
\]
Claim~\ref{c:shell-recolor} of Lemma~\ref{lem:shell} and Lemma~\ref{lem:word-reindexing}
show that $\Phi:\mathcal C_0(U)\to\mathcal C_0(V_*)$ is injective.
Its fibers therefore have size at most $1=\operatorname{Fib}_m(\nu,0)$.

Second, suppose $P\ge1$. Let us define $(Q_s,h_s,\eta_s)$ for $s\in[0,\ell]_{\Z}$.
Let $(Q_0,h_0,\eta_0)=(Y,h,\eta)$.
For $s\in[\ell]_{\Z}$, let us construct $(Q_s,h_s,\eta_s)$ recursively.
The run-edit $\omega_s$ acts on the copy of $U_{s-1}$ in $Q_{s-1}$.
Here $h_s=h_s(g,\gamma)$ and $\eta_s=\eta_s(g,\gamma)$.
Define $a_t\in[M]_{\Z}$ and $b_t\in[M]_{\Z}$ for $t\in[P]_{\Z}$ by
\[
\begin{aligned}
 W_{\mathrm L}&=\operatorname{Trans}(R_{a_P}\ccn \cdots\ccn R_{a_1},\operatorname{Start}_{\mathrm{pt}}(W_{\mathrm L})),\\
 W_{\mathrm R}&=\operatorname{Trans}(R_{b_1}\ccn \cdots\ccn R_{b_P},\operatorname{Start}_{\mathrm{pt}}(W_{\mathrm R})).
\end{aligned}
\]

Let us insert a simple run at step $s\in I_R$.
Let $Q_{s-1}=B_s^{\mathrm L}\ccn U_{s-1}\ccn B_s^{\mathrm R}$ and put $q_s=|B_s^{\mathrm L}|+i_s$.
The paths $B_s^{\mathrm L}$ and $B_s^{\mathrm R}$ contain the outer simple runs. Thus,
$q_s\in[1,|Q_{s-1}|-1]_{\Z}$.
Assumption~\ref{hyp:comparison-run-lengths} gives $j_s\in[M]_{\Z}\subseteq[m]_{\Z}$.
Apply Lemma~\ref{lem:simple-run-insertion} with
$Q=Q_{s-1}$, $q=q_s$, $j=j_s$, and $c=\operatorname{Surp}(Q_{s-1},h_{s-1})$.
Let $(h_s,\eta_s)$ be the image of $(h_{s-1},\eta_{s-1})$ under its map $\Phi_{q,j,c}$.
The resulting path is
\[
 Q_s=\operatorname{Ins}(Q_{s-1},q_s,R_{j_s})=B_s^{\mathrm L}\ccn U_s\ccn B_s^{\mathrm R}.
\]

Let us verify the hypotheses for the vertical run-edits.
At step $s\in I_+\cup I_-$, let
\[
 Q_{s-1}=B_s^{\mathrm L}\ccn U_{s-1}\ccn B_s^{\mathrm R},\qquad
 t_s=\#\bigl((I_+\cup I_-)\cap[s]_{\Z}\bigr).
\]
The paths $B_s^{\mathrm L}$ and $B_s^{\mathrm R}$ consist of the simple runs added around $U_{s-1}$.
Every simple run of $Q_{s-1}$ has vertical length in $[M]_{\Z}$
by assumptions~\ref{hyp:comparison-run-lengths}, \ref{hyp:comparison-shell-budget}, and~\ref{hyp:comparison-target}.

First, let us insert a vertical edge when $s\in I_+$.
Then, $A>0$, and assumption~\ref{hyp:comparison-insertion-range} gives
\[
 m\ge4,\qquad M\le m-2.
\]
Claims~\ref{c:vertical-insert-connected}--\ref{c:vertical-insert-injective}
of Lemma~\ref{lem:vertical-insert} therefore apply after the first $|B_s^{\mathrm L}|+i_s$ edges.

Second, let us delete a vertical edge when $s\in I_-$.
Let $v_s=(U_{s-1})_{r_s}$ and let
\[
 \widehat v_s=(Q_{s-1})_{|B_s^{\mathrm L}|+r_s}.
\]
Claims~\ref{c:vertical-delete-connected}--\ref{c:vertical-delete-map}
of Lemma~\ref{lem:vertical-delete} apply at $\widehat v_s$ because $M\le m-1$.
For each connected colored grading $(h,\eta)$ on $Q_{s-1}$, fix an edge
$y_s(h,\eta)$ satisfying Claim~\ref{c:vertical-delete-connected} of Lemma~\ref{lem:vertical-delete}
for $X=Q_{s-1}$, $v=\widehat v_s$, $u_0=a_{t_s}$, and $v_0=b_{t_s}$.
At step $s$, use $y_s(h_{s-1},\eta_{s-1})$.

Let us specify the shell path used at step $s\in I_+\cup I_-$.
Use the parameters $a_{t_s}$ and $b_{t_s}$ in the construction
\ref{hyp:insertion-shell}--\ref{hyp:insertion-colors} of Lemma~\ref{lem:vertical-insert}
or the construction \ref{hyp:deletion-shell}--\ref{hyp:deletion-color} of Lemma~\ref{lem:vertical-delete}, respectively.
The resulting path is
\[
 Q_s=R_{a_{t_s}}\ccn B_s^{\mathrm L}\ccn U_s\ccn B_s^{\mathrm R}\ccn R_{b_{t_s}}.
\]
Claim~\ref{c:vertical-insert-injective} of Lemma~\ref{lem:vertical-insert}
and Claim~\ref{c:vertical-delete-map} of Lemma~\ref{lem:vertical-delete}
supply the colored grading $(h_s,\eta_s)$.

Let us compute $\operatorname{Surp}(Q_\ell,h_\ell)$.
For $s\in[0,\ell]_{\Z}$, let $c_s=\operatorname{Surp}(Q_s,h_s)$.
The initial shell gives $c_0=0$.
Equation~\eqref{eq:run-insertion-surplus}, Claim~\ref{c:vertical-insert-residual} of Lemma~\ref{lem:vertical-insert},
and Claim~\ref{c:vertical-delete-residual} of Lemma~\ref{lem:vertical-delete} give
\[
 c_s-c_{s-1}=
 \begin{cases}
 1,&s\in I_R,\\
 0,&s\in I_+,\\
 h_{s-1}(\widehat v_s),&s\in I_-.
 \end{cases}
\]
Since $h_{s-1}(\widehat v_s)\in[0,m]_{\Z}$, summing gives
\[
 c_\ell=K+\sum_{s\in I_-}h_{s-1}(\widehat v_s)\le K+mD.
\]

Let us reduce the surplus to zero.
Let $n=P-A-D$.
Assumption~\ref{hyp:comparison-shell-budget} gives
\[
 n\ge K+mD\ge c_\ell.
\]
Apply Lemma~\ref{lem:surplus-reduction} with $Q=Q_\ell$ and
\[
 \alpha'=(a_{A+D+j})_{j\in[n]_{\Z}},\qquad
 \beta'=(b_{A+D+j})_{j\in[n]_{\Z}}.
\]
Let $Q'$ denote the path $Q_n$ in that application and let
\[
 (f,\delta)=\Delta_{Q_\ell,\alpha',\beta'}(h_\ell,\eta_\ell)
 \in\mathcal C_0(Q').
\]

Let us identify $Q'$ and define $\Phi$.
The paths $Q_\ell$ and $Q'$ are translated copies of
\[
 R_{a_{A+D}}\ccn \cdots\ccn R_{a_1}\ccn R_u\ccn V\ccn R_v\ccn R_{b_1}\ccn \cdots\ccn R_{b_{A+D}}
\]
and
\[
 R_{a_P}\ccn \cdots\ccn R_{a_1}\ccn R_u\ccn V\ccn R_v\ccn R_{b_1}\ccn \cdots\ccn R_{b_P},
\]
respectively. When $A+D=0$, the blocks indexed by $[A+D]_{\Z}$ are empty.
Let
\[
 \mathbf t=\operatorname{Start}_{\mathrm{pt}}(V_*)-
             \operatorname{Start}_{\mathrm{pt}}(Q').
\]
Then, $\operatorname{Trans}(Q',\mathbf t)=V_*$.
Using the translated grading and color function from
Definition~\ref{def:colored-grading}(5), define
\[
 \Phi(g,\gamma)=(f+\mathbf t,\delta+\mathbf t).
\]
Lemma~\ref{lem:word-reindexing} gives $\Phi(g,\gamma)\in\mathcal C_0(V_*)$.

Let us bound the fibers of $\Phi$ by recovering the colored gradings on $U$.
Let $D_U\subseteq N(U)$ be the set of edges of $U$ deleted by $\Omega$ in assumption~\ref{hyp:comparison-edit-list}.
For $e\in\operatorname{NE}(U)\setminus D_U$, let $\widetilde e\in\operatorname{NE}(V_*)$ be its image under the inclusion $U\to Q_0$, the edge maps $Q_{s-1}\to Q_s$ at steps $s=1,\ldots,\ell$, the inclusion $Q_\ell\to Q'$, and the translation by $\mathbf t$, in that order.
Fix $(h,\eta)\in\mathcal C_0(V_*)$.
Define $\widehat h\in G(U)$ by
\[
 \widehat h(e)=
 \begin{cases}
 h(\widetilde e),& e\notin D_U,\\
 0,& e\in D_U.
 \end{cases}
\]

Let us recover an input colored grading from its differences and deleted edges.
Let $e_1,\ldots,e_\nu$ be the horizontal edges of $U$ in path order.
For $(g,\gamma)\in\Phi^{-1}(\{(h,\eta)\})$, let
\[
 \delta_i=\operatorname{Diff}_U(g,\widehat h;e_i)
 =g(e_i)-h(\widetilde e_i)\qquad(i\in[\nu]_{\Z}).
\]
Claim~\ref{c:shell-recolor} of Lemma~\ref{lem:shell}, Claim~\ref{c:vertical-insert-preserved} of Lemma~\ref{lem:vertical-insert}, and Lemma~\ref{lem:surplus-reduction} preserve the grades and colors of the traced input edges.
Equation~\eqref{eq:run-insertion-grades} and Definition~\ref{def:recolor} give grade and color preservation for simple-run insertion.
Claim~\ref{c:vertical-delete-preserved} of Lemma~\ref{lem:vertical-delete} decreases one horizontal grade by $1$ and preserves all other surviving grades and colors.
Assumption~\ref{hyp:comparison-edit-list} specifies $D$ deletions.
Hence, $\delta_i\ge0$ for $i\in[\nu]_{\Z}$ and $\sum_{i=1}^{\nu}\delta_i\le D$.
If $\delta_i=0$, no deletion decreases the grade of the traced edge, so its color is preserved.
We recover $(g,\gamma)$ from:
\begin{enumerate}[label=(\roman*)]
\item the differences $\delta_1,\ldots,\delta_\nu$;
\item the colors $\gamma(e_i)$ for $i\in[\nu]_{\Z}$ such that $\delta_i>0$;
\item the pairs $(g(v),\gamma(v))$ for $v\in D_U$.
\end{enumerate}
We have
\[
\begin{aligned}
 g(e_i)&=h(\widetilde e_i)+\delta_i &&(i\in[\nu]_{\Z}),\\
 \gamma(e_i)&=\eta(\widetilde e_i) &&(\delta_i=0),\\
 g(v)&=h(\widetilde v),\qquad \gamma(v)=\eta(\widetilde v)
 &&(v\in N(U)\setminus D_U).
\end{aligned}
\]

Let us count the difference vectors, horizontal colors, and grade--color pairs in (i)--(iii).
The vector $(\delta_1,\ldots,\delta_\nu)$ has
$\binom{\nu+D}{D}$ possibilities. At most $D$
original horizontal colors change. Each has at most $2^m$ choices.
We delete at most $D$ original vertical edges. Each has
$\sum_{r=0}^m\binom mr=2^m$ possible grade--color pairs.
For each $\xi\in\mathcal C_0(V_*)$, we obtain
\[
 \#\Phi^{-1}(\{\xi\})\le\binom{\nu+D}{D}(2^m)^D(2^m)^D=\operatorname{Fib}_m(\nu,D).
\]

\prf{Claim~\ref{c:composition-bound}}
By Claim~\ref{c:composition-target},
\[
 C_0(U)=\sum_{\xi\in\mathcal C_0(V_*)}\#\Phi^{-1}(\{\xi\})
 \le \operatorname{Fib}_m(\nu,D)C_0(V_*).
\]
\end{proof}

\begin{definition}\label{def:entropy-loss-function}
For $x>0$, define
\[
 \psi_m(x)=(1+x)\log(1+x)-x\log x+2mx\log2.
\]
Let
\[
 \psi_m(0)=0.
\]
\end{definition}

\begin{lemma}\label{lem:entropy-loss}
Let $\nu\in\Z_{\ge0}$, $D\in\Z_{\ge0}$, $N\in\R_{>0}$, and $x\in\R_{\ge0}$ satisfy $\nu\le N$ and $D\le xN$.
Then,
\begin{equation}\label{eq:entropy-loss}
 \log \operatorname{Fib}_m(\nu,D)\le N\psi_m(x).
\end{equation}
\end{lemma}
\begin{proof}
First, suppose $x=0$. Then, $D=0$ and $\log \operatorname{Fib}_m(\nu,0)=0=N\psi_m(0)$.

Second, suppose $x>0$. Let us prove~\eqref{eq:entropy-loss}.
Let $t=\frac{x}{1+x}\in(0,1)$.
The binomial expansion gives
\[
 \binom{\nu+D}{D}t^D(1-t)^\nu
 \le\bigl(t+(1-t)\bigr)^{\nu+D}=1.
\]
Since $-\log t>0$ and $-\log(1-t)>0$, the bounds $\nu\le N$ and $D\le xN$ give
\[
\begin{aligned}
 \log \operatorname{Fib}_m(\nu,D)
 &\le-D\log t-\nu\log(1-t)+2mD\log2\\
 &\le-xN\log t-N\log(1-t)+2mxN\log2\\
 &=N\psi_m(x).
\end{aligned}
\]
\end{proof}

